\clearpage
\fancyhead[L]{\small HELLER FLECK IM MOND.}
\SourceLeaf{441}
\begin{center}
{\large [HELLER FLECK IM MOND.]}
\end{center}

\begin{center}
\noindent\rule{\linewidth}{0.4pt}

Göttingische gelehrte Anzeigen. Stück 46. Seite 449..452. 1821. März 22.

\noindent\rule{\linewidth}{0.4pt}
\end{center}

Ueber eine von Hrn. Dr. \textsc{Olbers} am 5 Febr. d. J. am dunkeln Theile der
Mondoberfläche beobachtete Erscheinung theilen wir hier einen Auszug aus einem Briefe
desselben an Hrn. Hofrath \textsc{Gauss}, vom 27. Februar, um so lieber mit, da die von
jenem genialen Astronomen beigefügten höchst sinnreichen Vermuthungen über das an sich
schon so seltene als merkwürdige Phänomen die aufmerksamste Beachtung verdienen und zu
einer ganz neuen und fruchtbaren Ansicht führen.

Am 5. Februar (schreibt Hr. Dr. \textsc{Olbers}) habe ich die Erscheinung im Monde
gesehen, die man einen Mondsvulkan genannt hat. Was ich in meinem Tagebuche darüber
sogleich niedergeschrieben habe, ist wörtlich folgendes: `Am fünften war es sehr heiter,
aber schon Mondschein. In dem dunkeln Theile des Mondes sah ich noch nie das Phänomen,
das man für einen brennenden Vulkan im Monde gehalten hat, so deutlich und auffallend,
wie diesen Abend. Es schien, wie gewöhnlich, im Aristarch zu sein. Es war klein, aber
ganz \emph{auffallend} heller, als der übrige Theil des von der Sonne nicht erleuchteten
Mondes, ganz sternähnlich, und hatte eben das Ansehen, wie ein Nordost vom Monde
stehender Fixstern 3ter Grösse.'

Da es am 6. Februar trübe war, hat Hr. Dr. \textsc{Olbers} sich nicht weiter nach dieser
Erscheinung umgesehen. Inzwischen haben, bald nachher, Englische öffentliche Blätter
angezeigt, dass der Capitain \textsc{Kater} am 7. Februar der Königl. Societät zu London
eine Nachricht über einen von ihm im Monde gesehenen Vulkan mitgetheilt habe. Er habe
sich durch fortgesetzte Beobachtungen wirklich überzeugt, dass es ein im Ausbruch
begriffener Vulkan sei.

Es scheint also (fährt Hr. Dr. \textsc{Olbers} fort), dass Hr. \textsc{Kater} dieselbe
Erscheinung gesehen, nur sie weiter verfolgt habe, die auch mir am 5. Februar auffiel.
Ich kenne zwar seine Ueberzeugungsgründe, dass dies wirklich ein brennender Vulkan
gewesen, nicht: allein nach allem, was wir von der Beschaffenheit des Mondes, und seiner
so zweifelhaften Atmosphäre wissen, scheint ein brennender Vulkan fast unmöglich.
`Vielmehr glaube ich, dass sich die Erde in einer ebenen merklich glatten fast einer
polirten Fläche ähnlichen Seitenwand einer zum Aristarch gehörenden grossen Felsklippe
wirklich abspiegelte.' Das so abgespiegelte Bild eines Theiles der Erde musste ganz
ungleich heller sein, als alles übrige bloss von der Erde erleuchtete, da dieses das
Erdenlicht nach allen Richtungen zerstreut, jenes dasselbe nur in \emph{einer} Richtung
zurückwirft. Wenn jene unvollkommene Spiegelung auch nur ein Zehntel des Erdlichts
zurückwarf (da unsre wirklichen Spiegel etwa die Hälfte des auf sie fallenden Lichtes
zurückwerfen) und die Seitenwand nur $2''$ im Durchmesser hatte, so konnte sie immer so
hell wie ein Stern sechster Grösse erscheinen. --- Nach dieser Vorstellung wird es
erklärlich: \emph{erstens}, warum wir die vulkanartigen Erscheinungen immer nur an
bestimmten Stellen des Mondes sehen. \emph{Zweitens}, warum sie nicht in jeder
Lunation, sondern nur selten zu Gesichte kommen: die Libration muss nemlich bis auf etwa
$2^\circ$ dieselbe sein. --- Die Möglichkeit, dass es solche mehr oder weniger
spiegelartig das Licht zurückwerfende Seitenwände der Mondsklippen geben könne, lässt
sich wohl nicht bezweifeln. Auf unsrer Erde mag es grosse Gletscherflächen geben, die
auch als unvollkommene Spiegel Licht zurückwerfen können. Ich führe dieses nur als etwas
analoges an: denn Gletscher sind im Monde eben so
\SourceLeaf{442}%
unwahrscheinlich, wie brennende Vulkane. Aber dass es unter den auch im Monde wahrscheinlich nach
Crystallisationsgesetzen gebildeten Gebirgen einzelne geben könne, die ebene, glatte, fast einer Politur
ähnliche Seitenflächen haben, scheint sehr gedenkbar. Es mag ihrer vielleicht viele im Monde geben,
aber selten mögen sie gerade die Lage haben, dass sie uns unter bestimmter Libration gerade das Bild
der Erde zurückspiegeln können, auch die Sonne scheint sich zuweilen auf ähnlichen Klippenwänden
im Monde abzuspiegeln. Noch vor etwa 8 Wochen sah ich im Mare Imbrium ausser der Lichtgrenze
zwei von der Sonne beschienene Bergköpfe, so ungewöhnlich hell, scintillirend, Siriusähnlich, dass es
mir unmöglich schien, hier bloss nach gewöhnlichen Zerstreuungsgesetzen zurückgeworfenes Sonnenlicht
anzunehmen. Ich kann nicht bestimmt sagen, ob der eine dieser Berge vielleicht der auch im Mare
Imbrium gelegene Lahire (nach \textsc{Schröter}) war, bei welchem \textsc{Schröter} ganz ähnliche Erscheinungen
wahrgenommen hat.

In einem spätern Briefe bemerkt Hr. Dr. \textsc{Olbers} noch, dass diese Hypothese über die Ursache
des Phänomens sich leicht prüfen lassen werde, weil, wenn sie die wahre Erklärung enthalte, bei derselben Libration immer dieselbe Erscheinung wieder statt haben müsse. Felsenwände, die ein Bild der
Erde oder der Sonne, mehr oder weniger unvollkommen zurückspiegeln können, seien im Monde um
so gedenkbarer, da dort wahrscheinlich nicht wie auf der Erde eine Verwitterung der äussern Oberfläche der Gebirge und Klippen durch atmosphärische Einwirkung statt finde. Die zurückspiegelnde
Klippenwand brauche auch nicht ganz eben zu sein, wenn sich nur die zurückspiegelnden Theile in parallelen Ebenen befinden, wie dies bei solchen Bergen, die nach Crystallisationsgesetzen gebildet sind,
leicht statt finden könne. Hr. Dr. \textsc{Olbers} erinnert hiebei an unsere Basaltberge, deren einzelne grosse
Crystalle noch sehr wohl dem entfernten Auge vereint ein unvollkommenes Sonnenbild zurückspiegeln
könnten, wenn ihre Oberflächen nicht längst durch Luft, Dünste, Regen u.\,s.\,w. die wahrscheinlich ursprünglich vorhanden gewesene Politur und Glätte verloren hätten.

Am 6. März, wo die Nachtseite des Mondes vortrefflich zu sehen war, konnte Hr. Dr. \textsc{Olbers}
mit seinem \textsc{Dollond}schen Fernrohr, alle Flecken, z.\,B. Grimaldi, Copernicus, Kepler, Maeilius, Menelaus u.\,s.\,w. sehr deutlich erkennen. Aristarch zeichnete sich wieder vor allen andern, auch, wie es
schien, \emph{mehr als gewöhnlich}, aus. Allein so \emph{hell} und so \emph{fixsternähnlich} wie am 5. Februar fand ihn Hr.
Dr. \textsc{Olbers} diesmal nicht.
